\documentclass[12pt]{article}

\usepackage[a4paper,margin=1in]{geometry}
\usepackage{amsmath,amssymb}
\usepackage{enumitem}

\setlength{\parindent}{0pt}
\setlength{\parskip}{0.7em}

\begin{document}

\begin{center}
    {\Large \textbf{Tutorial: Hitchhiker Problem}}\\
    \vspace{0.2cm}
    {\large Converting Arrival Time Questions into Poisson Count Questions}
\end{center}

\section*{Question}

Cars arrive at a rate of 10 per hour. Each car picks up a hitchhiker with probability
\[
\frac{1}{10}.
\]
You are second in line. What is the probability that you wait for more than 2 hours?

\section*{Main Confusing Step}

The step we want to understand is
\[
P(S_2>2)=P(N(2)<2).
\]

This file explains why that conversion is true.

\section*{Theory 1: Poisson Thinning}

Suppose cars arrive according to a Poisson process with rate
\[
\lambda = 10 \text{ cars per hour}.
\]

But not every car is useful to the hitchhiker. A car is useful only if it picks up a
hitchhiker.

Each car picks up a hitchhiker with probability
\[
p=\frac{1}{10}.
\]

By the thinning property of a Poisson process, the useful cars also form a Poisson
process, with rate
\[
\lambda_{\text{useful}} = \lambda p.
\]

Therefore,
\[
\lambda_{\text{useful}}
=
10\left(\frac{1}{10}\right)
=
1.
\]

So useful cars arrive at rate
\[
\boxed{1 \text{ useful car per hour}.}
\]

\section*{Theory 2: Meaning of \(N(t)\)}

Let
\[
N(t)=\text{number of useful cars that arrive by time }t.
\]

Since useful cars arrive according to a Poisson process with rate 1 per hour,
\[
N(t)\sim \text{Poisson}(t).
\]

For example:
\[
N(2)=\text{number of useful cars that arrive in the first 2 hours}.
\]

Because the rate is 1 per hour, the mean number of useful cars in 2 hours is
\[
1\times 2=2.
\]

So
\[
\boxed{N(2)\sim \text{Poisson}(2).}
\]

\section*{Theory 3: Meaning of \(S_2\)}

Let
\[
S_2=\text{time at which the second useful car arrives}.
\]

Since you are second in line, you need the second useful car.

\begin{itemize}
    \item The first useful car picks up the person in front of you.
    \item The second useful car picks you up.
\end{itemize}

So your waiting time is
\[
S_2.
\]

The question asks:
\[
P(S_2>2).
\]

This means:

\[
\text{Probability that the second useful car arrives after 2 hours.}
\]

\section*{Why \(P(S_2>2)=P(N(2)<2)\)}

The event
\[
S_2>2
\]
means:

\[
\text{The second useful car has not arrived by time 2.}
\]

If the second useful car has not arrived by time 2, then by time 2 there could be:
\[
0 \text{ useful cars}
\]
or
\[
1 \text{ useful car}.
\]

But there cannot be 2 or more useful cars by time 2.

So
\[
S_2>2
\]
is exactly the same event as
\[
N(2)<2.
\]

Therefore,
\[
\boxed{P(S_2>2)=P(N(2)<2).}
\]

\section*{Timeline Interpretation}

Think of a time line from 0 to 2 hours.

\[
\text{0 hours} \qquad \rule{6cm}{0.4pt} \qquad \text{2 hours}
\]

If by the end of this interval fewer than 2 useful cars arrived, then you are still
waiting after 2 hours.

That is why the problem changes from an arrival-time question into a count question.

\section*{Step-by-Step Calculation}

We already found that
\[
N(2)\sim \text{Poisson}(2).
\]

We need
\[
P(S_2>2).
\]

Convert it:
\[
P(S_2>2)=P(N(2)<2).
\]

The condition \(N(2)<2\) means
\[
N(2)=0 \quad \text{or} \quad N(2)=1.
\]

Therefore,
\[
P(N(2)<2)
=
P(N(2)=0)+P(N(2)=1).
\]

Use the Poisson formula:
\[
P(N=k)=e^{-\mu}\frac{\mu^k}{k!}.
\]

Here,
\[
\mu=2.
\]

So
\[
P(N(2)=0)
=
e^{-2}\frac{2^0}{0!}
=
e^{-2}.
\]

Also,
\[
P(N(2)=1)
=
e^{-2}\frac{2^1}{1!}
=
2e^{-2}.
\]

Therefore,
\[
P(S_2>2)
=
e^{-2}+2e^{-2}.
\]

Hence,
\[
\boxed{P(S_2>2)=3e^{-2}.}
\]

\section*{Final Answer}

\[
\boxed{P(\text{you wait more than 2 hours})=3e^{-2}.}
\]

Numerically,
\[
3e^{-2}\approx 0.406.
\]

So there is about a
\[
\boxed{40.6\%}
\]
chance that you wait more than 2 hours.

\section*{General Method}

For a Poisson process:

\[
S_n = \text{time of the }n\text{th event},
\]
and
\[
N(t)=\text{number of events by time }t.
\]

Then
\[
\boxed{S_n>t \iff N(t)<n.}
\]

Therefore,
\[
\boxed{P(S_n>t)=P(N(t)<n).}
\]

This becomes
\[
\boxed{
P(S_n>t)
=
\sum_{k=0}^{n-1} e^{-\lambda t}\frac{(\lambda t)^k}{k!}.
}
\]

\section*{Practice Question 1}

\textbf{Question.} Useful buses arrive according to a Poisson process with rate 3 per
hour. A student needs the third useful bus. What is the probability that the student
waits more than 1 hour?

\subsection*{Answer}

Let
\[
S_3=\text{arrival time of the third useful bus}.
\]

The question asks:
\[
P(S_3>1).
\]

Convert to a count:
\[
P(S_3>1)=P(N(1)<3).
\]

Since the rate is 3 per hour,
\[
N(1)\sim \text{Poisson}(3).
\]

Thus,
\[
P(N(1)<3)=P(N(1)=0)+P(N(1)=1)+P(N(1)=2).
\]

So
\[
P(S_3>1)
=
e^{-3}\frac{3^0}{0!}
+e^{-3}\frac{3^1}{1!}
+e^{-3}\frac{3^2}{2!}.
\]

Therefore,
\[
P(S_3>1)
=
e^{-3}\left(1+3+\frac{9}{2}\right).
\]

Hence,
\[
\boxed{P(S_3>1)=8.5e^{-3}.}
\]

\section*{Practice Question 2}

\textbf{Question.} Customers arrive according to a Poisson process with rate 5 per hour.
What is the probability that the fourth customer arrives within 30 minutes?

\subsection*{Answer}

Let
\[
S_4=\text{arrival time of the fourth customer}.
\]

We want
\[
P(S_4\leq 0.5).
\]

The event
\[
S_4\leq 0.5
\]
means the fourth customer has arrived by 0.5 hours.

So by 0.5 hours, there must be at least 4 customers:
\[
P(S_4\leq 0.5)=P(N(0.5)\geq 4).
\]

Since the rate is 5 per hour,
\[
N(0.5)\sim \text{Poisson}(5\times 0.5)=\text{Poisson}(2.5).
\]

Therefore,
\[
P(S_4\leq 0.5)=P(N(0.5)\geq 4).
\]

Use the complement:
\[
P(N(0.5)\geq 4)=1-P(N(0.5)\leq 3).
\]

So
\[
P(S_4\leq 0.5)
=
1-\sum_{k=0}^{3}e^{-2.5}\frac{2.5^k}{k!}.
\]

Thus,
\[
\boxed{
P(S_4\leq 0.5)
=
1-\left[
e^{-2.5}
+e^{-2.5}\frac{2.5}{1!}
+e^{-2.5}\frac{2.5^2}{2!}
+e^{-2.5}\frac{2.5^3}{3!}
\right].
}
\]

\section*{Key Exam Shortcut}

Use these conversions:

\[
\boxed{S_n>t \iff N(t)<n}
\]

\[
\boxed{S_n\leq t \iff N(t)\geq n}
\]

So:

\begin{center}
\begin{tabular}{|c|c|}
\hline
\textbf{Arrival-time statement} & \textbf{Poisson-count statement}\\
\hline
\(n\)th event happens after time \(t\) & fewer than \(n\) events by time \(t\)\\
\hline
\(n\)th event happens by time \(t\) & at least \(n\) events by time \(t\)\\
\hline
\end{tabular}
\end{center}

\end{document}
